%페이지 세팅
\documentclass[10pt,a4paper]{article}
\setlength{\parindent}{0em}                  %DISTANCIA SANGRÍA
\setlength{\parskip}{0.5em}                  %DISTANCIA ENTRE PÁRRAFOS
\textwidth 6.5in
\textheight 9.in
\oddsidemargin 0in
\headheight 0in

\usepackage{amsmath}
\usepackage{tcolorbox}
\usepackage{amssymb}
\usepackage{amsthm}
\usepackage{lastpage}
\usepackage{fancyhdr}
\usepackage{accents}
\usepackage{setup}
\usepackage{import}
\usepackage{fancyhdr}
\usepackage{layouts}
\addtolength{\voffset}{0mm}
\addtolength{\textheight}{0mm}

\usepackage{xcolor}
\usepackage{mdframed}
\usepackage[shortlabels]{enumitem}
\usepackage{indentfirst}
\usepackage{hyperref}

\makeatletter
%%%%%%%%%%%%%%%%%%%%%%%%%%%%%%%%%%%%%%%%%%%%%%%%%%%%%%%%%%%%%%%%%%%%%%%%%%
% Sectioning Sytle Change
\renewcommand{\thesubsection}{\thesection-\arabic{subsection}}
% \renewcommand{\thesubsection}{\thesection.\alph{subsection}}
% Customize Section Style
% we use \prefix@<level> only if it is defined
\renewcommand{\@seccntformat}[1]{%
  \ifcsname prefix@#1\endcsname
    \csname prefix@#1\endcsname
  \else
    \csname the#1\endcsname\quad
  \fi}
% define \prefix@section
\newcommand\prefix@section{Topic  \thesection\ : }
% \newcommand\prefix@subsection{Experiment \thesubsection\ :\ }
\makeatother




\renewcommand{\thesubsection}{\thesection\ - \arabic{subsection}}
\linespread{1.15}
%%%%%%%%%%%%%%%%%%%%%%%%%%%%%%%%%%%%%%%%%%%%%%%%%%%%%%%%%%%%%%%%%%%%%%%%%%
    % 매주 리포트를 작성할때 이 부분을 수정하면 보고서 전체가 수정된다. 
%%%%%%%%%%%%%%%%%%%%%%%%%%%%%%%%%%%%%%%%%%%%%%%%%%%%%%%%%%%%%%%%%%%%%%%%%%
% 작성하는 주차 
\newcommand{\numnum}{9}
% 이번 실험의 부제목 
% \newcommand{\subtitle}{wireshark(2) : TCP UDP IP protocols}
% \newcommand{\subtitle}{USPR(1) : Introduction, LABVIEW, Spectagram}
% \newcommand{\subtitle}{USPR(2) : Radar generator tutorial}
% \newcommand{\subtitle}{USPR(3) : MAC control, CNN protector}
% \newcommand{\subtitle}{Mininet(1) : SDN-based Switch and Hub}
% \newcommand{\subtitle}{Mininet(2) : SDN-based Routing}
\newcommand{\subtitle}{Socket Programming}
%-------------------------------------------------------------------------
%--> 떠다니는 객체 사용하는 부분 
\usepackage{wrapfig}
\newenvironment{problem}[2][Problem]  
    { \begin{mdframed}[backgroundcolor=gray!20] \textbf{#1 #2} \\}
    {  \end{mdframed}}
%-------------------------------------------------------------------------
\begin{document}
\pagestyle{fancy}
\fancyhf{}
    \rhead{\small 2조 2016142096 조윤신, 2017142043 김재민}
    \lhead{\textsc{Pre Report Week \numnum \quad \subtitle}}
\cfoot{\thepage}
\renewcommand\headrulewidth{0.3mm}
%\renewcommand\footrulewidth{0.3mm}
\thispagestyle{plain}
\begin{flushleft}
\textsc{School of Electrical and Electronic of Enginnering} \\
\textsc{eee 4474-01 : Experiments on Communication Networks }\\[0.1cm]
\small{\textsc{\textbf{2조} 2016142096 \textbf{조윤신}, 2017142043 \textbf{김재민}}}\\
\end{flushleft}
%-------------------------------------------------------------------------
\begin{flushright}\vspace{-25mm}
    \includegraphics[height=3cm]{image/1.jpg}
    \vspace{5mm}
\end{flushright}
    \begin{center}\vspace{-1.5cm}
        \textbf{\huge  Pre Report Week\numnum}\\ 
        \vspace{0.6mm}
        {\large \subtitle}\\           
    \end{center}
\vspace{-5mm}
\rule{\linewidth}{0.4mm}
\vspace{-5mm}
% main documents
% 모든 프로그래밍 편집기 (IDE)에서 원도우 기준으로 '컨트롤 + /'를 누르면 
% 커서가 활성화된 line의 주석처리를 해제 지정할 수 있다. 
%%%%%%%%%%%%%%%%%%%%%%%%%%%%%%%%%%%%%%%%%%%%%%%%%%%%%%%%%%%%%%%%%%%%%%%%%%
                        % For Week 01 ~ 09.14
%%%%%%%%%%%%%%%%%%%%%%%%%%%%%%%%%%%%%%%%%%%%%%%%%%%%%%%%%%%%%%%%%%%%%%%%%%
% \import{./TEX/week01}{01_01}
% \newpage
% \import{./TEX/week01}{01_02}
% \newpage
% \import{./TEX/week01}{01_03}
% \newpage
% \import{./TEX/week01}{02_udp}
% \newpage
% \import{./TEX/week01}{03_01}
% \newpage
% \import{./TEX/week01}{03_02}
% \import{./TEX/week01}{03_03}
%-------------------------------------------------------------------------
%%%%%%%%%%%%%%%%%%%%%%%%%%%%%%%%%%%%%%%%%%%%%%%%%%%%%%%%%%%%%%%%%%%%%%%%%%
                        % For Week 02 ~ 09.21
%%%%%%%%%%%%%%%%%%%%%%%%%%%%%%%%%%%%%%%%%%%%%%%%%%%%%%%%%%%%%%%%%%%%%%%%%%
% \import{./TEX/week02}{01_01}
% \import{./TEX/week02}{02_01}
%-------------------------------------------------------------------------
%%%%%%%%%%%%%%%%%%%%%%%%%%%%%%%%%%%%%%%%%%%%%%%%%%%%%%%%%%%%%%%%%%%%%%%%%%
                        % For Week 03 ~ 09.28
%%%%%%%%%%%%%%%%%%%%%%%%%%%%%%%%%%%%%%%%%%%%%%%%%%%%%%%%%%%%%%%%%%%%%%%%%%
% \import{./TEX/week03}{01_chirp}
% \import{./TEX/week03}{02_5g}
% \import{./TEX/week03}{03_}
%-------------------------------------------------------------------------
%%%%%%%%%%%%%%%%%%%%%%%%%%%%%%%%%%%%%%%%%%%%%%%%%%%%%%%%%%%%%%%%%%%%%%%%%%
                        % For Week 04 ~ 10.05
%%%%%%%%%%%%%%%%%%%%%%%%%%%%%%%%%%%%%%%%%%%%%%%%%%%%%%%%%%%%%%%%%%%%%%%%%%
% \import{./TEX/week04}{01}
% \import{./TEX/week04}{02}
% \clearpage
% \import{./TEX/week04}{03}
% \import{./TEX/week04}{04}
%-------------------------------------------------------------------------
%%%%%%%%%%%%%%%%%%%%%%%%%%%%%%%%%%%%%%%%%%%%%%%%%%%%%%%%%%%%%%%%%%%%%%%%%%
                        % For Week 05 ~ 10.12
%%%%%%%%%%%%%%%%%%%%%%%%%%%%%%%%%%%%%%%%%%%%%%%%%%%%%%%%%%%%%%%%%%%%%%%%%%
% \import{./TEX/week05}{01}
% \import{./TEX/week05}{02}
% \import{./TEX/week05}{03}
% \import{./TEX/week05}{04}
%-------------------------------------------------------------------------
%%%%%%%%%%%%%%%%%%%%%%%%%%%%%%%%%%%%%%%%%%%%%%%%%%%%%%%%%%%%%%%%%%%%%%%%%%
                        % For Week 06 ~ 10.19
%%%%%%%%%%%%%%%%%%%%%%%%%%%%%%%%%%%%%%%%%%%%%%%%%%%%%%%%%%%%%%%%%%%%%%%%%%
% \import{./TEX/week06}{01}
% \import{./TEX/week06}{02}
% \import{./TEX/week06}{03}
% \import{./TEX/week06}{04}
% \import{./TEX/week06}{05}
%-------------------------------------------------------------------------
%%%%%%%%%%%%%%%%%%%%%%%%%%%%%%%%%%%%%%%%%%%%%%%%%%%%%%%%%%%%%%%%%%%%%%%%%%
                        % For Week 07 ~ 10.26
%%%%%%%%%%%%%%%%%%%%%%%%%%%%%%%%%%%%%%%%%%%%%%%%%%%%%%%%%%%%%%%%%%%%%%%%%%
\import{./TEX/week07}{01}
\import{./TEX/week07}{02}

% \import{./TEX/week03}{01_}
% \import{./TEX/week03}{02_}
% \import{./TEX/week03}{03_}
%-------------------------------------------------------------------------
\end{document}